% Generated by manual/tools/gen_reference.py from the code. Do not edit by hand.

\begin{reftable}{tandem-dash launch}{tab:ref-launch-flags}{\hd{Flag} & \hd{Default} & \hd{What it does}}
\texttt{ENGINE} & \texttt{} & the engine to start: vllm or sglang. \\
\texttt{-{}-model} & \texttt{required} & the model id or path the engine serves. \\
\texttt{-{}-port} & \texttt{none} & the engine's port. \\
\texttt{-{}-host} & \texttt{127.0.0.1} & the engine's bind address. \\
\texttt{-{}-id} & \texttt{none} & the backend id (default: the engine's name and port, e.g. vllm8100). \\
\texttt{-{}-label} & \texttt{none} & the backend's label in the dashboard. \\
\texttt{-{}-config-\allowbreak{}db} & \texttt{\textasciitilde{}/\allowbreak{}.tandemllm/\allowbreak{}config.sqlite3 (TANDEM\_\allowbreak{}CONFIG\_\allowbreak{}DB)} & the proxy's config database the backend is registered in. \\
\texttt{-{}-collect-\allowbreak{}url} & \texttt{none (TANDEM\_\allowbreak{}COLLECT\_\allowbreak{}BASE)} & the proxy address the plugin posts to; default: this machine's proxy on its port. \\
\texttt{-{}-engine-\allowbreak{}cmd} & \texttt{none} & the engine's command instead of \texttt{vllm serve} / \texttt{sglang serve}. \\
\texttt{-{}-python} & \texttt{none} & the Python of the engine's environment, for the plugin check (default: the one behind the engine's command). \\
\texttt{-{}-no-\allowbreak{}register} & \texttt{off (a switch)} & do not create or update the backend. \\
\texttt{-{}-dry-\allowbreak{}run} & \texttt{off (a switch)} & print, start nothing. \\
\end{reftable}
